\documentclass[11pt]{article}
\usepackage[paperwidth=4.5in,paperheight=6in,margin=0.625in]{geometry}
\usepackage{fontspec}
\usepackage{adjustbox}
\usepackage{xspace}

\usepackage{fancyhdr}
\pagestyle{fancy}
\fancyhf{}
\renewcommand{\headrulewidth}{0pt}
\fancyfoot[C]{\thepage}

\usepackage{setspace}
\usepackage{titlesec}
\titleformat{\section}{\normalfont\Large\bfseries\centering}{}{0pt}{}
\titleformat{\subsection}{\normalfont\large\bfseries\centering}{}{0pt}{}

\usepackage{multicol}
\setlength{\columnsep}{1cm}

\usepackage{tikz}
\usepackage{pgfornament}
\usetikzlibrary{calc}
\usepackage{varwidth}

\usepackage{polyglossia}
\setmainlanguage{hebrew}
\setotherlanguage{english}

%\newfontfamily\hebrewfont[Script=Hebrew,Scale=1.25]{Taamey Frank CLM}
\newfontfamily\hebrewfont[Script=Hebrew,Scale=1.25]{Drugulin CLM}
%\newfontfamily\titlefont[Script=Hebrew,Scale=1.25]{Yiddishkeit AlefAlefAlef Bold}
\newfontfamily\titlefont[Script=Hebrew,Scale=1.25]{Drugulin CLM}
\newfontfamily\mikra[Script=Hebrew,Scale=1.2]{Keter Aram Tsova}
\newfontfamily\rashi[Script=Hebrew,Scale=1.25]{Noto Rashi Hebrew}
\newfontfamily\plain[Script=Hebrew,Scale=1.25]{David Libre Medium}

\newcommand{\instruction}[1]{%
  \noindent{\plain\tiny #1}
}

% formats

\newcommand{\wholeline}[1]{%
\\
\adjustbox{max width=\textwidth}
{\RL{
\fontsize{100}{120}
#1
}}
}

\newcommand{\mekor}[1]{%
  \marginpar{\tiny #1}
}

% abbrevs

\newcommand{\ps}[0]{תה׳\xspace}

\newfontfamily\englishfont{TeX Gyre Pagella}

% pretty

% corner side corner_size side_len textwidth
\newenvironment{ornatebox}[5]{%
  \def\corner{#1}
  \def\side{#2}
  \def\size{#3}
  \def\sidelen{#4}
  \centering
    \begin{tikzpicture}[every node/.style={inner sep=0pt}]
      % open the node and the varwidth so that user content goes inside it
      \node[inner sep=2mm, align=justify, text width=#5] (box)
      \bgroup%
      \RL \bgroup%
      \setlength{\parindent}{0pt}
      \setlength{\leftskip}{0pt plus -0.5fil}%
      \setlength{\rightskip}{0pt plus 0.5fil}%
      %\parfillskip 0pt plus fil%
        %\begin{varwidth}{\textwidth}\centering\ignorespaces
}{% <-- end of begin-code, start of end-code
        % close varwidth, close node (the semicolon ends the \node)
        %\unskip\end{varwidth}
      \par
      \egroup \egroup; % end \node

      \node[anchor=center](CNW)
      at (box.north west) {\pgfornament[width=\size]{\corner}};
      \node[anchor=center](CNE)
      at (box.north east) {\pgfornament[width=\size,symmetry=v]{\corner}};
      \node[anchor=center](CSW)
      at (box.south west) {\pgfornament[width=\size,symmetry=h]{\corner}};
      \node[anchor=center](CSE)
      at (box.south east) {\pgfornament[width=\size,symmetry=c]{\corner}};

      % measure box size
      \path let \p1 = (box.north east), \p2 = (box.south west) in
        \pgfextra{
          \pgfmathsetmacro{\boxwidth}{\x1-\x2-2*\size}
          \pgfmathsetmacro{\boxheight}{\y1-\y2-2*\size}
          \pgfmathsetmacro{\boxwidthin}{\boxwidth/\sidelen}
          \pgfmathsetmacro{\boxheightin}{\boxheight/\sidelen}
          \pgfmathtruncatemacro{\numh}{max(1,ceil(\boxwidthin))}
          \pgfmathtruncatemacro{\numv}{max(1,ceil(\boxheightin))}
          \pgfornamentline{CNE.north west}{CNW.north east}{\numh}{\side}
          \pgfornamentline{CSW.south east}{CSE.south west}{\numh}{\side}
          \pgfornamentline{CNW.south west}{CSW.north west}{\numv}{\side}
          \pgfornamentline{CSE.north east}{CNE.south east}{\numv}{\side}
        };

    \end{tikzpicture}
    \par
}

\newenvironment{floralbox}[1][0.6\textwidth]{%
  \begin{ornatebox}{63}{70}{0.6in}{0.75in}{#1}
  %\begin{ornatebox}{39}{46}{0.4in}{3}{1}
}{
  \end{ornatebox}
}

\newenvironment{simplebox}[1][0.8\textwidth]{%
  \begin{ornatebox}{37}{85}{0.3in}{1in}{#1}
  %\begin{ornatebox}{39}{46}{0.4in}{3}{1}
}{
  \end{ornatebox}
}

% RTL
\usepackage{bidi}

% no numbers in toc
\setcounter{secnumdepth}{0}

\usepackage{sectsty}
\allsectionsfont{\titlefont}

\tolerance=2000
\emergencystretch=1.5em
\setlength{\parskip}{0.5em}
\interlinepenalty=10000 % keep pars together
\setlength{\parindent}{0pt}

\begin{document}
\pagenumbering{gobble}

% Title page

\newgeometry{left=0pt,right=0pt,bottom=0.25in}

{
\centering
\vspace{0.5in}
{\plain בענטשער}\\
\vspace{0.125in}
\begin{floralbox}
%\centering
\fontsize{40}{60}\selectfont
טוֹבִ֥ים\; דֹּדֶ֖יךָ \par מִיָּֽיִן%
\end{floralbox}
%\wholeline{\textbf{טוֹבִ֥ים דֹּדֶ֖יךָ מִיָּֽיִן}}
\vspace{0.5in}
{\plain שמחת הנישואין של} \\
\vspace{0.125in}
{\Huge דוד אריה וולך} \\
{\Huge וסימן־טוב רזיאל קליינר} \\
\vspace{0.125in}
{\plain ראש חודש טבת תשפ״ו}

}
\clearpage
\restoregeometry

\tableofcontents
\newpage

\section{ברכת המזון}

\begin{samepage}
\subsection{תהלים לפני זימון}

\instruction{ביום שאומרים בו תחנון:}

עַל\mekor{\ps קלז}נַהֲרוֹת בָּבֶל שָׁם יָשַֽׁבְנוּ גַּם בָּכִֽינוּ בְּזׇכְרֵֽנוּ אֶת צִיּוֹן׃ עַל
עֲרָבִים בְּתוֹכָהּ תָּלִֽינוּ כִּנֹּרוֹתֵֽינוּ׃ כִּי שָׁם שְׁאֵלֽוּנוּ שׁוֹבֵֽינוּ דִּבְרֵי שִׁיר וְתוֹלָלֵֽינוּ שִׂמְחָה
שִׁירוּ לָֽנוּ מִשִּׁיר צִיּוֹן׃ אֵיךְ נָשִׁיר אֶת שִׁיר יְהוָה עַל אַדְמַת נֵכָר׃ אִם אֶשְׁכָּחֵךְ
יְרוּשָׁלָ‍ִם תִּשְׁכַּח יְמִינִי׃ תִּדְבַּק לְשׁוֹנִי לְחִכִּי אִם לֹא אֶזְכְּרֵֽכִי אִם לֹא אַעֲלֶה אֶת
יְרוּשָׁלַ‍ִם עַל רֹאשׁ שִׂמְחָתִי׃ זְכֹר יְהוָה לִבְנֵי אֱדוֹם אֵת יוֹם יְרוּשָׁלָ‍ִם
הָאֹמְרִים עָֽרוּ עָֽרוּ עַד הַיְסוֹד בָּהּ׃ בַּת בָּבֶל הַשְּׁדוּדָה אַשְׁרֵי שֶׁיְשַׁלֶּם לָךְ אֶת גְּמוּלֵךְ שֶׁגָּמַלְתְּ
לָֽנוּ׃ אַשְׁרֵי שֶׁיֹּאחֵז וְנִפֵּץ אֶת עֹלָלַֽיִךְ אֶל הַסָּֽלַע׃ 
\end{samepage}

\begin{samepage}
\instruction{ביום שאין אומרים בו תחנון:}

שִׁיר\mekor{\ps קכו}הַמַּעֲלוֹת בְּשׁוּב יְהוָה אֶת שִׁיבַת צִיּוֹן הָיִֽינוּ כְּחֹלְמִים׃ אָז
יִמָּלֵא שְׂחוֹק פִּינוּ וּלְשׁוֹנֵֽנוּ רִנָּה אָז יֹאמְרוּ בַגּוֹיִם הִגְדִּיל יְהוָה לַעֲשׂוֹת עִם אֵֽלֶּה׃
הִגְדִּיל יְהוָה לַעֲשׂוֹת עִמָּֽנוּ הָיִֽינוּ שְׂמֵחִים׃ שׁוּבָה יְהוָה אֶת שְׁבִיתֵֽנוּ כַּאֲפִיקִים בַּנֶּֽגֶב׃
הַזֹּרְעִים בְּדִמְעָה בְּרִנָּה יִקְצֹֽרוּ׃ הָלוֹךְ יֵלֵךְ וּבָכֹה נֹשֵׂא מֶֽשֶׁךְ הַזָּֽרַע בֹּא יָבוֹא בְרִנָּה נֹשֵׂא
אֲלֻמֹּתָיו׃ 
\end{samepage}

\newpage

\begin{samepage}
\subsection{זימון}

\instruction{מזמן:}\\רבותי, מיר וועלן בענטשן \instruction{(י״א: רַבּוֹתַי, נְבָרֵךְ)}

\instruction{מסבים:}\\{\mikra יְהִ֤י\mekor{\ps קיג ב} שֵׁ֣ם יְהֹוָ֣ה מְבֹרָ֑ךְ מֵ֝עַתָּ֗ה וְעַד־עוֹלָֽם:}

\instruction{מזמן:}\\בִּרְשׁוּת מָרָנָן וְרַבָּנָן וְרַבּוֹתַי,\\נְבָרֵךְ (אֱלֹהֵינוּ) שֶׁאָכַלְנוּ מִשֶּׁלּוֹ.

\instruction{מסבים:}\\ בָּרוּךְ (אֱלֹהֵינוּ) שֶׁאָכַלְנוּ מִשֶּׁלּוֹ וּבְטוּבוֹ חָיִֽינוּ.

\instruction{מזמן:}\\ בָּרוּךְ (אֱלֹהֵינוּ) שֶׁאָכַלְנוּ מִשֶּׁלּוֹ וּבְטוּבוֹ חָיִֽינוּ.
\end{samepage}

\newpage

\begin{samepage}
\subsection*{ברכת הזן}
בָּרוּךְ אַתָּה יְיָ אֱלֹהֵֽינוּ מֶֽלֶךְ הָעוֹלָם, הַזָּן אֶת הָעוֹלָם כֻּלּוֹ בְּטוּבוֹ בְּחֵן בְּחֶֽסֶד וּבְרַחֲמִים,
הוּא {\mikra נֹתֵ֣ן\mekor{\ps קלו כה}לֶ֭חֶם לְכׇל־בָּשָׂ֑ר כִּ֖י לְעוֹלָ֣ם חַסְדּֽוֹ}. וּבְטוּבוֹ
הַגָּדוֹל תָּמִיד לֹא חָסַר לָֽנוּ וְאַל יֶחְסַר לָֽנוּ מָזוֹן לְעוֹלָם וָעֶד, בַּעֲבוּר שְׁמוֹ הַגָּדוֹל, כִּי
הוּא אֵל זָן וּמְפַרְנֵס לַכֹּל וּמֵטִיב לַכֹּל, וּמֵכִין מָזוֹן לְכׇל בְּרִיּוֹתָיו אֲשֶׁר בָּרָא. כָּאָמוּר,
\mekor{\ps קמה טז}
\wholeline{\mikra{פּוֹתֵ֥חַ אֶת־יָדֶ֑ךָ וּמַשְׂבִּ֖יעַ לְכׇל־חַ֣י רָצֽוֹן:}}\par
{\Large\centering {בָּרוּךְ אַתָּה יְיָ הַזָּן אֶת הַכֹּל.}\par}
\end{samepage}

\begin{samepage}
\subsection*{ברכת הארץ}
נוֹדֶה לְךָ יְיָ אֱלֹהֵֽינוּ, עַל שֶׁהִנְחַֽלְתָּ לַאֲבוֹתֵֽינוּ אֶֽרֶץ חֶמְדָּה טוֹבָה וּרְחָבָה, וְעַל
שֶׁהוֹצֵאתָֽנוּ יְיָ אֱלֹהֵֽינוּ מֵאֶֽרֶץ מִצְרַֽיִם וּפְדִיתָֽנוּ מִבֵּית עֲבָדִים, וְעַל בְּרִיתְךָ שֶׁחָתַֽמְתָּ
בִּבְשָׂרֵֽנוּ וְעַל תּוֹרָתְךָ שֶׁלִּמַּדְתָּֽנוּ, וְעַל חֻקֶּֽיךָ שֶׁהוֹדַעְתָּֽנוּ, וְעַל חַיִּים חֵן וָחֶֽסֶד
שֶׁחוֹנַנְתָּֽנוּ, וְעַל אֲכִילַת מָזוֹן שָׁאַתָּה זָן וּמְפַרְנֵס אוֹתָֽנוּ תָּמִיד, בְּכׇל יוֹם וּבְכׇל עֵת
וּבְכׇל שָׁעָה.
\end{samepage}

\begin{samepage}
\instruction{בחנוכה ובפורים:}

\begin{center}
{\Large עַל־הַנִּסִּים וְעַל־הַפֻּרְקָן וְעַל־הַגְּבוּרוֹת וְעַל־הַתְּשׁוּעוֹת וְעַל־הַמִּלְחָמוֹת שֶׁעָשִׂיתָ
לַאֲבוֹתֵינוּ בַּיָּמִים הָהֵם בַּזְּמַן הַזֶּה}
\end{center}
\end{samepage}

\begin{samepage}
\begin{simplebox}
\instruction{בחנוכה:}\\
בִּימֵי מַתִּתְיָהוּ בֶן יוֹחָנָן כֹּהֵן גָּדוֹל חַשְׁמוֹנָאִי וּבָנָיו כְּשֶׁעָמְדָה מַלְכוּת יָוָן הָרְשָׁעָה
עַל־עַמְּךָ יִשְׂרָאֵל לְהַשְׁכִּיחָם תּוֹרָתֶךָ וּלְהַעֲבִירָם מֵחֻקֵּי רְצוֹנֶךָ. וְאַתָּה, בְּרַחֲמֶיךָ הָרַבִּים
עָמַדְתָּ לָהֶם בְּעֵת צָרָתָם, רַבְתָּ אֶת־רִיבָם, דַּנְתָּ אֶת־דִּינָם, נָקַמְתָּ אֶת־נִקְמָתָם. מָסַרְתָּ
גִבּוֹרִים בְּיַד חַלָּשִׁים וְרַבִּים בְּיַד מְעַטִּים וּטְמֵאִים בְּיַד טְהוֹרִים וּרְשָׁעִים בְּיַד צַדִּיקִים
וְזֵדִים בְּיַד עוֹסְקֵי תוֹרָתֶךָ. וּלְךָ עָשִׂיתָ שֵׁם גָּדוֹל וְקָדוֹשׁ בְּעוֹלָמֶךָ, וּלְעַמְּךָ יִשְׂרָאֵל
עָשִׂיתָ תְּשׁוּעָה גְדוֹלָה וּפֻרְקָן כְּהַיּוֹם הַזֶּה. וְאַחַר־כָּךְ בָּאוּ בָנֶיךָ לִדְבִיר בֵּיתֶךָ, וּפִנּוּ
אֶת־הֵיכָלֶךָ וְטִהֲרוּ אֶת־מִקְדָּשֶׁךָ וְהִדְלִיקוּ נֵרוֹת בְּחַצְרוֹת קָדְשֶׁךָ, וְקָבְעוּ שְׁמוֹנַת יְמֵי
חֲנֻכָּה אֵלּוּ לְהוֹדוֹת וּלְהַלֵּל לְשִׁמְךָ הַגָּדוֹל.
\end{simplebox}

\begin{simplebox}
\instruction{בפורים:}\\
בִּימֵי מָרְדְּכַי וְאֶסְתֵּר בְּשׁוּשַׁן הַבִּירָה, כְּשֶׁעָמַד עֲלֵיהֶם הָמָן הָרָשָׁע, {\mikra
לְהַשְׁמִ֡יד\mekor{אסתר ג יג}לַהֲרֹ֣ג וּלְאַבֵּ֣ד אֶת־כׇּל־הַ֠יְּהוּדִ֠ים מִנַּ֨עַר וְעַד־זָקֵ֜ן טַ֤ף
וְנָשִׁים֙ בְּי֣וֹם אֶחָ֔ד בִּשְׁלוֹשָׁ֥ה עָשָׂ֛ר לְחֹ֥דֶשׁ שְׁנֵים־עָשָׂ֖ר הוּא־חֹ֣דֶשׁ אֲדָ֑ר וּשְׁלָלָ֖ם לָבֽוֹז:}
וְאַתָּה בְּרַחֲמֶיךָ הָרַבִּים הֵפַרְתָּ אֶת־עֲצָתוֹ וְקִלְקַלְתָּ אֶת־מַחֲשַׁבְתּוֹ וַהֲשֵׁבֹתָ לוֹ גְּמוּלוֹ
בְרֹאשׁוֹ, וְתָלוּ אוֹתוֹ וְאֶת בָּנָיו עַל הָעֵץ.
\end{simplebox}
\end{samepage}

\begin{samepage}
וְעַל הַכֹּל יְיָ אֱלֹהֵֽינוּ אֲנַֽחְנוּ מוֹדִים לָךְ וּמְבָרְכִים אוֹתָךְ, יִתְבָּרַךְ שִׁמְךָ בְּפִי כׇּל חַי
תָּמִיד לְעוֹלָם וָעֶד, כַּכָּתוּב:
\mekor{דברים ח י}\wholeline{\mikra וְאָכַלְתָּ֖ וְשָׂבָ֑עְתָּ וּבֵֽרַכְתָּ֙ אֶת־יְהֹוָ֣ה אֱלֹהֶ֔יךָ}\\
\wholeline{\mikra עַל־הָאָ֥רֶץ הַטֹּבָ֖ה אֲשֶׁ֥ר נָֽתַן־לָֽךְ:}\par
{\Large\centering בָּרוּךְ אַתָּה יְיָ עַל הָאָֽרֶץ וְעַל הַמָּזוֹן.\par}
\end{samepage}

\newpage

\begin{samepage}
\subsection*{ברכת בונה ירושלים}
רַחֵם יְיָ אֱלֹהֵֽינוּ עַל יִשְׂרָאֵל עַמֶּֽךָ, וְעַל יְרוּשָׁלַֽיִם עִירֶֽךָ, וְעַל צִיּוֹן מִשְׁכַּן כְּבוֹדֶֽךָ,
וְעַל מַלְכוּת בֵּית דָּוִד מְשִׁיחֶֽךָ, וְעַל הַבַּֽיִת הַגָּדוֹל וְהַקָּדוֹשׁ שֶׁנִּקְרָא שִׁמְךָ עָלָיו. אֱלֹהֵֽינוּ,
אָבִֽינוּ, רְעֵֽנוּ, זוּנֵֽנוּ, פַּרְנְסֵֽנוּ וְכַלְכְּלֵֽנוּ וְהַרְוִיחֵֽנוּ, וְהַרְוַח לָֽנוּ יְיָ אֱלֹהֵֽינוּ מְהֵרָה
מִכׇּל צָרוֹתֵֽינוּ. וְנָא אַל תַּצְרִיכֵֽנוּ יְיָ אֱלֹהֵֽינוּ, לֹא לִידֵי מַתְּנַת בָּשָׂר וָדָם וְלֹא לִידֵי
הַלְוָאָתָם, כִּי אִם לְיָדְךָ הַמְּלֵאָה הַפְּתוּחָה הַקְּדוֹשָׁה וְהָרְחָבָה, שֶׁלֹּא נֵבוֹשׁ וְלֹא נִכָּלֵם לְעוֹלָם
וָעֶד.
\end{samepage}

\begin{samepage}
\instruction{בשבת:}\\
רְצֵה וְהַחֲלִיצֵֽנוּ יְיָ אֱלֹהֵֽינוּ בְּמִצְוֹתֶֽיךָ וּבְמִצְוַת יוֹם הַשְּׁבִיעִי הַשַּׁבָּת הַגָּדוֹל וְהַקָּדוֹשׂ
הַזֶּה. כִּי יוֹם זֶה גָּדוֹל וְקָדוֹשׁ הוּא לְפָנֶֽיךָ לִשְׁבׇּת בּוֹ וְלָנוּחַ בּוֹ בְּאַהֲבָה כְּמִצְוַת
רְצוֹנֶֽךָ. וּבִרְצוֹנְךָ הָנִֽיחַ לָֽנוּ יְיָ אֱלֹהֵֽינוּ שֶׁלֹּא תְהֵא צָרָה וְיָגוֹן וַאֲנָחָה בְּיוֹם
מְנוּחָתֵֽנוּ. וְהַרְאֵֽנוּ יְיָ אֱלֹהֵֽינוּ בְּנֶחָמַת צִיוֹן עִירֶֽךָ וּבְבִנְיַן יְרוּשָׁלַֽיִם עִיר קׇדְשֶֽׁךָ כִּי
אַתָּה הוּא בַּֽעַל הַיְשׁוּעוֹת וּבַֽעַל הַנֶּחָמוֹת.
\end{samepage}

\begin{samepage}
\instruction{בראש חודש ובמועדים:}\\
אֱלֹהֵינוּ וֵאלֹהֵי אֲבוֹתֵינוּ, יַעֲלֶה וְיָבֹא וְיַגִּיעַ, וְיֵרָאֶה וְיֵרָצֶה וְיִשָּׁמַע, וְיִפָּקֵד וְיִזָּכֵר,
זִכְרוֹנֵנוּ וּפִקְדוֹנֵנוּ וְזִכְרוֹן אֲבוֹתֵינוּ, וְזִכְרוֹן מָשִׁיחַ בֶּן דָּוִד עַבְדֶּךָ, וְזִכְרוֹן
יְרוּשָׁלַיִם עִיר קׇדְשֶׁךָ, וְזִכְרוֹן כׇּל עַמְּךָ בֵּית יִשְׂרָאֵל, לְפָנֶיךָ, לִפְלֵטָה לְטוֹבָה לְחֵן לְחֶסֶד
וּלְרַחֲמִים, לְחַיִּים וּלְשָׁלוֹם, בְּיוֹם
\\\instruction{בראש חודש:}  רֹאשׁ הַחֹדֶשׁ הַזֶּה.
\\\instruction{בראש השנה:}  הַזִּכָּרוֹן הַזֶּה.
\\\instruction{בפסח:}  חַג הַמַּצּוֹת הַזֶּה.
\\\instruction{בשבועות:}  חַג הַשָּׁבוּעוֹת הַזֶּה.
\\\instruction{בסוכות:}  חַג הַסֻּכּוֹת הַזֶּה.
\\\instruction{בשמיני עצרת:}  הַשְּׁמִינִי חַג הָעֲצֶרֶת הַזֶּה.
\\זׇכְרֵנוּ יְיָ אֱלֹהֵינוּ בּוֹ לְטוֹבָה, וּפׇקְדֵנוּ בוֹ לִבְרָכָה, וְהוֹשִׁיעֵנוּ בוֹ לְחַיִּים. וּבִדְבַר
יְשׁוּעָה וְרַחֲמִים חוּס וְחׇנֵּנוּ, וְרַחֵם עָלֵינוּ וְהוֹשִׁיעֵנוּ, כִּי אֵלֶיךָ עֵינֵינוּ, כִּי אֵל מֶלֶךְ
חַנּוּן וְרַחוּם אָתָּה.
\end{samepage}

\begin{samepage}
וּבְנֵה יְרוּשָׁלַֽיִם עִיר הַקֹּֽדֶשׁ בִּמְהֵרָה בְיָמֵֽינוּ. בָּרוּךְ אַתָּה יְיָ, בּוֹנֵה בְרַחֲמָיו יְרוּשָׁלָֽיִם. אָמֵן.
\end{samepage}

\newpage

\begin{samepage}
\subsection*{ברכת הטוב והמיטיב}
בָּרוּךְ אַתָּה יְיָ אֱלֹהֵֽינוּ, מֶֽלֶךְ הָעוֹלָם, הָאֵל אָבִֽינוּ, מַלְכֵּֽנוּ, אַדִּירֵֽנוּ, בּוֹרְאֵֽנוּ,
גֹּאֲלֵֽנוּ, יוֹצְרֵֽנוּ, קְדוֹשֵֽׁנוּ קְדוֹשׁ יַעֲקֹב, רוֹעֵֽנוּ רוֹעֵה יִשְׂרָאֵל, הַמֶּֽלֶךְ הַטּוֹב וְהַמֵּיטִיב
לַכֹּל, שֶׁבְּכׇל יוֹם וָיוֹם הוּא הֵיטִיב, הוּא מֵיטִיב, הוּא יֵיטִיב לָֽנוּ, הוּא גְמָלָֽנוּ, הוּא
גוֹמְלֵֽנוּ, הוּא יִגְמְלֵֽנוּ לָעַד, לְחֵן וּלְחֶֽסֶד וּלְרַחֲמִים וּלְרֶֽוַח הַצָּלָה וְהַצְלָחָה, בְּרָכָה
וִישׁוּעָה, נֶחָמָה פַּרְנָסָה וְכַלְכָּלָה וְרַחֲמִים וְחַיִּים וְשָׁלוֹם, וְכׇל טוֹב; וּמִכׇּל טוּב לְעוֹלָם
אַל יְחַסְּרֵֽנוּ.
\end{samepage}

\subsection*{הרחמן}
הָרַחֲמָן הוּא יִמְלֹךְ עָלֵֽינוּ לְעוֹלָם וָעֶד.

הָרַחֲמָן הוּא יִתְבָּרַךְ בַּשָּׁמַיִם וּבָאָֽרֶץ.

הָרַחֲמָן הוּא יִשְׁתַּבַּח לְדוֹר דּוֹרִים, וְיִתְפָּאַר בָּֽנוּ לָעַד וּלְנֵצַח נְצָחִים, וְיִתְהַדַּר בָּֽנוּ לָעַד וּלְעוֹלְמֵי עוֹלָמִים.

הָרַחֲמָן הוּא יְפַרְנְסֵֽנוּ בְּכָבוֹד.

הָרַחֲמָן הוּא יִשְׁבֹּר עֹל גָּלֻיּוֹת מֵעַל צַוָּארֵֽנוּ, וְהוּא יוֹלִיכֵֽנוּ קוֹמְמִיוּת לְאַרְצֵֽנוּ.

הָרַחֲמָן הוּא יִשְׁלַח לָֽנוּ בְּרָכָה מְרֻבָּה בַּבַּיִת הַזֶּה, וְעַל שֻׁלְחָן זֶה שֶׁאָכַֽלְנוּ עָלָיו (שֻׁלְחָנוֹת אֵֽלּוּ שֶׁאָכַֽלְנוּ עֲלֵיהֶם).

הָרַחֲמָן הוּא יִשְׁלַח לָֽנוּ אֶת אֵלִיָּֽהוּ הַנָּבִיא זָכוּר לַטּוֹב, וִיבַשֶּׂר לָֽנוּ בְּשׂוֹרוֹת טוֹבוֹת יְשׁוּעוֹת וְנֶחָמוֹת.

ברכת האורח
יְהִי רָצוֹן, שֶׁלֹּא יֵבוֹשׁ בַּֽעַל הַבַּֽיִת בָּעוֹלָם הַזֶּה, וְלֹא יִכָּלֵם לָעוֹלָם הַבָּא, וְיִצְלַח מְאֹד בְּכׇל נְכָסָיו, וְיִהְיוּ נְכָסָיו וּנְכָסֵֽינוּ מֻצְלָחִים וּקְרוֹבִים לָעִיר, וְאַל יִשְׁלֹט שָׂטָן לֹא בְּמַעֲשֵׂי יָדָיו וְלֹא בְּמַעֲשֵׂי יָדֵֽינוּ, וְאַל יִזְדַקֵּק (נוסח הגמרא: יִזְדַקֵּר) לֹא לְפָנָיו וְלֹא לְפָנֵֽינוּ שׁוּם דְבַר הִרְהוּר חֵטְא וַעֲבֵרָה וְעָוֹן מֵעַתָּה וְעַד עוֹלָם.

הָרַחֲמָן הוּא יְבָרֵךְ אֶת (אָבִי מוֹרִי) בַּֽעַל הַבַּֽיִת הַזֶּה, וְאֶת (אִמִּי מוֹרָתִי) בַּעֲלַת הַבַּֽיִת הַזֶּה, אוֹתָם וְאֶת בֵּיתָם וְאֶת זַרְעָם וְאֶת כׇּל אֲשֶׁר לָהֶם, אוֹתָֽנוּ וְאֶת כׇּל אֲשֶׁר לָֽנוּ, כְּמוֹ שֶׁנִּתְבָּרְכוּ אֲבוֹתֵֽינוּ אַבְרָהָם יִצְחָק וְיַעֲקֹב "בַּכֹּל" (בראשית כד א)־"מִכֹּל" (בראשית כז לג)־"כֹּל" (בראשית לג יא) – כֵּן יְבָרֵךְ אוֹתָֽנוּ כֻּלָּֽנוּ יַֽחַד בִּבְרָכָה שְׁלֵמָה. וְנֹאמַר: "אָמֵן".

על שולחן עצמו אומר:
הָרַחֲמָן הוּא יְבָרֵךְ אוֹתִי, (אם יש לו אב: וְאֶת אָבִי מוֹרִי,) (אם יש לו אם: וְאֶת אִמִּי מוֹרָתִי,) (נשוי: וְאֶת אִשְׁתִּי,) (נשואה: וְאֶת בַּעְלִי,) וְאֶת זַרְעִי, וְאֶת כׇּל אֲשֶׁר לִי, אוֹתָֽנוּ וְאֶת כׇּל אֲשֶׁר לָֽנוּ, כְּמוֹ שֶׁנִּתְבָּרְכוּ אֲבוֹתֵֽינוּ אַבְרָהָם יִצְחָק וְיַעֲקֹב "בַּכֹּל" (בראשית כד א)־"מִכֹּל" (בראשית כז לג)־"כֹּל" (בראשית לג יא) – כֵּן יְבָרֵךְ אוֹתָֽנוּ כֻּלָּֽנוּ יַֽחַד בִּבְרָכָה שְׁלֵמָה. וְנֹאמַר: "אָמֵן".
בסעודה משותפת: הָרַחֲמָן הוּא יְבָרֵךְ אֶת כׇּל הַמְסֻבִּין כַּאן. אוֹתָם וְאֶת בֵּיתָם וְאֶת זַרְעָם וְאֶת כׇּל אֲשֶׁר לָהֶם.

בַּמָּרוֹם יְלַמְּדוּ (עֲלֵיהֶם וְ)עָלֵֽינוּ זְכוּת שֶׁתְּהֵא לְמִשְׁמֶֽרֶת שָׁלוֹם. וְנִשָּׂא בְרָכָה מֵאֵת יְיָ, וּצְדָקָה מֵאֱלֹהֵי יִשְׁעֵֽנוּ, וְנִמְצָא חֵן וְשֵֽׂכֶל טוֹב בְּעֵינֵי אֱלֹהִים וְאָדָם.

בשבת, חג או ראש חודש[הראה]
יש מוסיפים:
הָרַחֲמָן הוּא יְבָרֵךְ אֶת מְדִינַת יִשְׂרָאֵל, רֵאשִׁית צְמִיחַת גְּאֻלָּתֵֽנוּ.

הָרַחֲמָן הוּא יְבָרֵךְ אֶת חַיָּלֵי צְבָא הַהֲגַנָּה לְיִשְׂרָאֵל וְאַנְשֵׁי כֹּחוֹת הַבִּטָּחוֹן, הָעוֹמְדִים עַל מִשְׁמַר אַרְצֵֽנוּ וְעָרֵי אֱלֹהֵֽינוּ, מִגְּבוּל הַלְּבָנוֹן וְעַד מִדְבַּר מִצְרַֽיִם, וּמִן הַיָּם הַגָּדוֹל עַד לְבוֹא הָעֲרָבָה, בַּיַּבָּשָׁה בָּאֲוִיר וּבַיָּם.

הָרַחֲמָן הוּא יְזַכֵּֽנוּ לִימוֹת הַמָּשִֽׁיחַ וּלְחַיֵּי הָעוֹלָם הַבָּא.

מַגְדִּיל [ביום שיש בו מוסף ובמלווה מלכה: מִגְדּוֹל] יְשׁוּעוֹת מַלְכּוֹ, וְעֹֽשֶׂה חֶֽסֶד לִמְשִׁיחוֹ, לְדָוִד וּלְזַרְעוֹ עַד עוֹלָם: (\ps יח נא, ש"ב כב נא)

עֹשֶׂה שָׁלוֹם בִּמְרוֹמָיו (איוב כה ב) הוּא יַעֲשֶׂה שָלוֹם עָלֵֽינוּ וְעַל כׇּל יִשְׂרָאֵל וְאִמְרוּ: "אָמֵן".

יְראוּ אֶת יְיָ קְדֹשָׁיו כִּי אֵין מַחְסוֹר לִירֵאָיו׃ כְּפִירִים רָשׁוּ וְרָעֵבוּ וְדֹרְשֵׁי יְיָ לֹא יַחְסְרוּ כׇל טוֹב׃ (\ps לד י-יא)
הוֹדוּ לַייָ כִּי טוֹב כִּי לְעוֹלָם חַסְדּוֹ׃ (\ps קו א, \ps קז א, \ps קיח א, \ps קיח כט, \ps קלו א, דברי הימים א טז לד)
פּוֹתֵֽחַ אֶת יָדֶֽךָ וּמַשְׂבִּֽיעַ לְכׇל חַי רָצוֹן׃ (\ps קמה טז)
בָּרוּךְ הַגֶּֽבֶר אֲשֶׁר יִבְטַח בַּייָ וְהָיָה יְיָ מִבְטַחוֹ׃ (ירמיהו יז ז)
נַֽעַר הָיִֽיתִי גַּם זָקַֽנְתִּי וְלֹא רָאִֽיתִי צַדִּֽיק נֶעֱזָב וְזַרְעוֹ מְבַקֶּשׁ לָֽחֶם: (\ps לז כה)
יְיָ עֹז לְעַמּוֹ יִתֵּן יְיָ יְבָרֵךְ אֶת עַמּוֹ בַשָּׁלוֹם: (\ps כט יא)

\end{document}
